\documentclass[twocolumn]{aastex631}
\begin{document}
\title{A residual reionization ionizing-photon-budget shortfall for star-forming galaxies at z\~{}6 (crisis systematic corner)}
\author{NebulaMind Lab (autonomous pipeline)}
\affiliation{NebulaMind Lab, \url{https://lab.nebulamind.net}}
\begin{abstract}
Reionization ionizing-photon-budget reconciliation at z\~{}6: star-forming galaxies require f\_esc=0.180 (+0.143/-0.079) to close the budget (Madau-Dickinson SFRD, log xi\_ion=25.5+/-0.15, clumping C=2-5, JWST-SFRD tail), versus indirect-proxy-inferred f\_esc=0.050 (+0.075/-0.030) from LzLCS O32/beta calibrations. Median delta(required-inferred)=+0.119 dex-frac (16-84\%: +0.020 to +0.264); 88\% of the systematic MC shows a shortfall. Conclusion: a genuine SHORTFALL remains, and the sign holds under both O32 and beta calibrations. The result is bounded by the xi\_ion x clumping x proxy-calibration systematic, not by statistics. Generated autonomously from public data (jwst) via the NebulaMind Lab runner: a bounded, reproducible, descriptive study.
\end{abstract}
\section{Introduction}
Recent studies have highlighted a potential crisis in the photon budget required for reionization [Muoz2024]. This has led to increased demands on ionizing sources and raised questions about the role of star-forming galaxies in this process [Davies2021, Park2022]. To address these concerns, it is essential to reconcile the ionizing-photon-budget using established literature values. Previous analyses have shown that understanding the galaxy ionizing photon budget is crucial for powering reionization [Duncan2015], and an analytic approach can provide valuable insights into this process [Madau2017].
\section{Data and method}
In this work, we adopt a literature-anchored budget calculation to assess the role of star-forming galaxies in reionization. We utilize the cosmic SFRD from Madau \& Dickinson (2014) and published values for xi\_ion and O32/beta f\_esc proxy calibrations [Chisholm+22, Flury+22; Simmonds+24]. Our method focuses on reconciling the ionizing-photon-budget without relying on new survey catalog data or observational results from specific telescopes.
\section{Result}
Our analysis reveals that star-forming galaxies require a escape fraction of f\_esc=0.180 (+0.143/-0.079) to close the reionization photon budget at z\~{}6, based on the Madau-Dickinson SFRD, log xi\_ion=25.5+/-0.15, and clumping factor C=2-5. However, indirect-proxy-inferred f\_esc from LzLCS O32/beta calibrations yields a value of 0.050 (+0.075/-0.030). This discrepancy results in a median delta(required-inferred)=+0.119 dex-frac (16-84\%: +0.020 to +0.264), with 88\% of the systematic Monte Carlo simulations showing a shortfall.
\begin{figure}\centering\includegraphics[width=\columnwidth]{result.png}\caption{A residual reionization ionizing-photon-budget shortfall for star-forming galaxies at z\~{}6 (crisis systematic corner)}\end{figure}
\section{Caveats}
It is essential to acknowledge that our approach has limitations, primarily due to its reliance on automated, single-selection, and uncalibrated measurements. The accuracy of our results depends heavily on the assumptions made in the literature values we adopt, such as the SFRD and xi\_ion calibrations. Additionally, our method does not account for potential variations in these parameters across different galaxy populations or redshifts. Furthermore, the use of proxy calibrations introduces uncertainty, as they may not accurately represent the true escape fraction. These limitations highlight the need for further research and more precise measurements to fully understand the reionization process.
\section*{References}
\footnotesize
[Muoz2024] Muñoz et al. (2024). Reionization after JWST: a photon budget crisis?. bibcode 2024MNRAS.535L..37M \par [Davies2021] Davies et al. (2021). The Predicament of Absorption-dominated Reionization: Increased Demands on Ionizing Sources. bibcode 2021ApJ...918L..35D \par [Park2022] Park et al. (2022). Calibrating excursion set reionization models to approximately conserve ionizing photons. bibcode 2022MNRAS.517..192P \par [Duncan2015] Duncan et al. (2015). Powering reionization: assessing the galaxy ionizing photon budget at z < 10. bibcode 2015MNRAS.451.2030D \par [Madau2017] Madau (2017). Cosmic Reionization after Planck and before JWST: An Analytic Approach. bibcode 2017ApJ...851...50M

\end{document}
